%==========================================================
\section{Reflexão sobre o uso de IA}\label{sec:reflexao}

Nas partes descritas nesta secção, a equipa idealizou primeiro cada componente à mão, a partir do relatório da Entrega~1.
Só depois usou a IA (Claude, Anthropic), com o mesmo ponto de partida, para comparar as duas versões. O objetivo não
foi substituir o trabalho da equipa, mas encontrar o que faltava em cada versão e juntar o melhor das duas, para chegar à
versão mais completa possível.

\subsection{\textit{Story map}}

Depois de idealizado o \textit{story map}, foram dados à IA o relatório da primeira entrega e os enunciados da primeira e da
segunda entregas, com o seguinte \textit{prompt}: «O ficheiro Elicitation é o enunciado da primeira entrega de um trabalho que
estou a fazer em grupo para uma cadeira de Engenharia de Software. O relatório que entregámos é o documento main.pdf. O
ficheiro Description é o enunciado da segunda entrega. A minha parte é fazer o story map. Consegues explicar-me o que é e
fazer-me um story board baseado no relatório que entregámos?»

A resposta da IA foi muito próxima do que procurávamos. Dividiu o mapa nas três \textit{releases} e propôs histórias de
utilizador curtas, escritas segundo os critérios INVEST, como pede o enunciado, mesmo sem lhe termos pedido as histórias.
Comparada com a versão da equipa, das 27 histórias propostas, 11 ficaram como estavam, 13 precisaram de pequenos ajustes e 3
tiveram de ser divididas em duas. Com a US-31, que saiu de um critério da US-15, ficaram 31 histórias.

\subsection{Histórias de utilizador}

As histórias foram depois revistas uma a uma, com a ajuda da IA, segundo os critérios INVEST e os dados da Entrega~1.
O \textit{prompt} foi apenas «review João's stories with the user-stories skill», sem dizer que as histórias tinham partido
de uma proposta da IA, para a revisão não ser enviesada. A revisão confirmou que os 61 registos e \textit{breakdowns} citados
existem, e apontou o que faltava: critérios de aceitação sem a forma \textit{Dado / quando / então} e sem caso de falha,
palavras de interface («Caderno», «quadro», «riscado»), histórias com duas capacidades e provas do agente usadas em histórias
do organizador. Também encontrou uma contradição: a US-16 dizia que o organizador «só reserva quando todos concordam», mas o
Roamly não faz reservas.

A IA propôs um plano de reescrita, mas as decisões foram da equipa: um só nome para o papel, «organizador(a) do grupo»;
histórias divididas com números novos (US-28 a US-31), sem sufixos a/b; e, em todas as histórias, a indicação de que a
reserva é feita fora do Roamly, no site de cada fornecedor. A equipa decidiu também passar a US-11 e a US-12 para a
Release~1, para esta cobrir a jornada toda. A IA concordou, mas desaconselhou mover também a US-07, porque depende de dados de
transportes que ninguém sabe ainda onde obter, e a US-07 ficou na Release~2. A numeração antiga do \textit{backlog} da
Entrega~1 deixou de ser usada, porque a IA não a conhecia e não acrescentava nada.

\subsection{Diagrama de casos de uso}

O diagrama de casos de uso foi pedido à IA a partir das histórias da Release~1, do material da cadeira (Cockburn, Gilson et
al., os \textit{slides} de UML) e de uma disposição definida pela equipa. A primeira versão não foi aceite: estava em inglês,
tinha subfunções que não tinham sido pedidas e não tinha pontos de extensão. Ao perguntarmos de onde vinham algumas regras,
percebemos que a IA apresentava decisões de desenho próprias como se fossem teoria e citava um critério que não está no material
da cadeira. A partir daí, demos-lhe os PDFs da cadeira e as decisões passaram a ser nossas: só as subfunções dadas como
exemplo, a IA do Roamly como parte do sistema, e não como ator, e os sites dos fornecedores como sistema externo. Corrigimos
também problemas visuais (linhas a cruzar elipses, texto cortado) e, depois da revisão do professor, pedimos a disposição
vertical, com o organizador ligado diretamente a cada objetivo.

\subsection{Casos de uso}

A reflexão sobre os casos de uso está na Secção~\ref{sec:uc-ia}, como pede o enunciado: a equipa reviu as primeiras versões,
com base em Cockburn e no UED, e a Tabela~\ref{tab:uc-ia} mostra o que mudou e porquê.

\subsection{Cenários BDD}

A equipa fixou primeiro as regras dos cenários, a partir de Dan North e do enunciado: um só comportamento por cenário, um só
\textit{Quando}, resultados observáveis, a linguagem dos casos de uso e etiquetas que ligam cada cenário à extensão, às
histórias e aos \textit{breakdowns}. Com essas regras e os casos de uso já corrigidos, a IA ajudou a escrever os primeiros
cenários, que a equipa reviu. Na revisão, retirámos um preço de voo que não vinha dos dados, passámos para um \textit{Dado} o
motivo que estava junto do acontecimento no \textit{Quando} e reescrevemos o cenário da extensão 3a do UC-04, que estava
confuso, para mostrar só o caso de um membro que não pode ir nas datas propostas. Como as entrevistas não dão datas exatas, as
datas de março usadas nos cenários são exemplos, iguais em todos os ficheiros.

\subsection{\textit{Domain-Driven Design}}

O modelo de domínio e os contextos delimitados também foram idealizados primeiro pela equipa. Depois, foram feitos dois
conjuntos de \textit{prompts} à IA, um para o modelo de domínio e outro para os contextos delimitados, e as respostas foram
comparadas com a versão da equipa.

\subsubsection{Modelo de domínio}

Para o modelo de domínio, a IA tinha como contexto o trabalho feito até ao \textit{story map}, e foram precisos três
\textit{prompts}. O primeiro foi: «My new task is now this: “Domain model (class-style diagram) in the ubiquitous language:
Trip, Proposal, Option (flight/stay/activity), Budget, Total per person, Group, Member/Coordinator, Vote/Opinion,
Availability, Decision, Price alert, Expense share, Payment… Derived from the artifact and sequence models”, can you do
it?» Como o \textit{prompt} estava em inglês, a resposta também veio em inglês, e o segundo \textit{prompt} serviu só para a
traduzir para português, por coerência com o resto do trabalho. O terceiro corrigiu incoerências entre a resposta e o
projeto, porque a IA ainda não conhecia as alterações que a equipa tinha feito às histórias.

Estes ajustes deveram-se ao contexto que demos à IA, e não à qualidade da resposta. A IA ajudou sobretudo nos diagramas,
que são a parte mais demorada de fazer à mão, e devolveu também o glossário que usou, sem lho termos pedido.

\subsubsection{Contextos delimitados}

Para os contextos delimitados, a IA tinha o mesmo contexto, já com as histórias atualizadas. O \textit{prompt} foi: «Tens
razão, as User stories foram alteradas pelo meu grupo e eu esqueci-me de tas dar. Aqui estão elas atualizadas. Agora, dou-te
a minha última tarefa: “At least two bounded contexts and a context map (e.g. Trip Planning \& Pricing / Group Decision /
Group Payments). Say what each owns, which terms change meaning across them (e.g. ‘Proposal’, ‘Member’), and the
relationships between them.” Quero a resposta em português».

O resultado ficou muito próximo da versão da equipa. A diferença principal foi o termo «coordenador» em vez de
«organizador». A IA disse tê-lo tirado do enunciado, mas veio do nosso próprio \textit{prompt}, que dizia «Member/Coordinator».
A lição é que a IA repete os termos que lhe damos, por isso o \textit{prompt} tem de usar a linguagem ubíqua do projeto.
